\documentclass[conference]{IEEEtran}
\usepackage{cite}
\usepackage{amsmath,amssymb,amsfonts}
\usepackage{graphicx}
\usepackage{textcomp}
\usepackage{xcolor}
\usepackage{booktabs}
\usepackage{hyperref}
\usepackage{listings}
\usepackage{algorithm}
\usepackage{algpseudocode}

\begin{document}

\title{PInsight: In-Situ Workload-Adaptive KV Cache Optimization for Long-Context LLM Inference}

\author{\IEEEauthorblockN{Joseph Bostok}
\IEEEauthorblockA{Department of Computer Science\\
University of North Carolina at Charlotte\\
Charlotte, NC, USA}}

\maketitle

% ============================================================
% ABSTRACT
% ============================================================
\begin{abstract}
The key-value (KV) cache is the central scalability bottleneck in long-context large language model (LLM) inference. Two families of optimization have emerged: \textbf{token eviction} (H2O, SnapKV), which scores and discards individual tokens, and \textbf{semantic chunking} (ChunkKV), which evicts contiguous token groups to preserve coherence. Both are evaluated with static configurations on standard benchmarks, yet the optimal strategy differs fundamentally across workload types---a chatbot benefits from aggressive sliding-window pruning, while retrieval-augmented generation (RAG) requires high-budget chunk-level eviction to preserve document integrity.

We present \textbf{PInsight}, an in-situ workload-adaptive optimization framework for vLLM that closes the loop between profiling and tuning. PInsight operates in four stages: (1)~\textit{profile} KV cache metrics during live serving, (2)~\textit{classify} the incoming workload by analyzing prompt structure and attention patterns, (3)~\textit{tune} cache parameters across three levels---token-level pruning, model-level chunking, and system-level vLLM configuration---matched to the detected workload, and (4)~\textit{monitor} the resulting performance to refine future decisions. We evaluate PInsight on four workload classes (dialogue, RAG, code generation, chain-of-thought reasoning) using Qwen2.5-7B-Instruct on NVIDIA H100 GPUs, comparing baseline vLLM against workload-tuned configurations. Our results demonstrate that adaptive tuning improves throughput by up to \textbf{XX\%} while reducing peak KV cache memory usage by \textbf{XX\%}, with the magnitude of improvement varying significantly across workload types---confirming that no single static configuration is optimal.
\end{abstract}

% ============================================================
% 1. INTRODUCTION
% ============================================================
\section{Introduction}

\subsection{The KV Cache Problem}

The deployment of large language models in production systems has exposed a fundamental tension between autoregressive generation and the memory-bandwidth constraints of modern GPU hardware. During decode, each new token requires reading the entire cached key-value state from GPU high-bandwidth memory (HBM)---a cost that grows linearly with sequence length. On an NVIDIA A100, single-token decode operates at approximately 1 FLOP/byte, a factor of $208\times$ below the GPU's compute ridge point~\cite{roofline}. The decode phase is firmly memory-bandwidth limited.

For Qwen2.5-7B with grouped-query attention (GQA)~\cite{gqa}, the KV cache requires approximately 57\,KB per token. Table~\ref{tab:kv_scaling} shows how this scales.

\begin{table}[h]
\centering
\caption{KV Cache Scaling for Qwen2.5-7B (BF16, 4 KV heads)}
\label{tab:kv_scaling}
\begin{tabular}{@{}rrc@{}}
\toprule
\textbf{Seq Length} & \textbf{KV Cache} & \textbf{\% of Weight BW} \\
\midrule
2,048 tokens & 117\,MB & 0.8\% \\
8,192 tokens & 470\,MB & 3.1\% \\
32,768 tokens & 1.9\,GB & 12.3\% \\
131,072 tokens & 7.5\,GB & 49.3\% \\
\bottomrule
\end{tabular}
\end{table}

\subsection{Different Workloads, Different Needs}

The KV cache problem manifests differently across workload types, each imposing distinct requirements on cache management:

\textbf{Multi-turn dialogue.} Conversational systems maintain full histories across many turns. Attention concentrates on initial ``sink'' tokens and the most recent turns, making aggressive pruning of middle turns safe. A 50-turn conversation reaches 10K+ context tokens.

\textbf{Retrieval-augmented generation (RAG).} RAG systems inject 5--20 retrieved passages into the prompt. These passages must remain \textit{intact}---per-token eviction fragments document context and degrades retrieval accuracy. Chunk-level eviction preserves passage integrity.

\textbf{Code generation.} Variable definitions, function signatures, and architectural constraints from hundreds of tokens earlier must remain available. Important tokens are \textit{scattered} across the sequence, requiring high retention budgets and fine-grained, token-level eviction.

\textbf{Chain-of-thought reasoning.} Long reasoning traces contain many redundant tokens (``let me reconsider,'' ``going back to step 1''). Aggressive pruning with redundancy detection is safe and significantly reduces memory pressure.

\textbf{The core insight:} a static KV cache configuration optimized for dialogue will \textit{hurt} RAG accuracy; one preserving RAG passages will \textit{waste memory} on dialogue workloads. No single configuration is optimal across workload types.

\subsection{The Optimization Landscape}

Two families of KV cache optimization address this problem:

\textbf{Token-level eviction.} H2O~\cite{h2o} retains ``heavy hitter'' tokens by cumulative attention score plus a sliding window. SnapKV~\cite{snapkv} uses per-head observation windows. StreamingLLM~\cite{streamingllm} preserves attention sinks plus a fixed window. These achieve up to $10\times$ compression but fragment semantic context.

\textbf{Semantic chunking.} ChunkKV~\cite{chunkkv} groups tokens into contiguous chunks and evicts at the chunk level, preserving coherence. It outperforms per-token methods by up to 8.7\% on long-context benchmarks but operates at coarser granularity.

Both approaches are evaluated with \textit{fixed} parameters on standard benchmarks. The choice between strategies---and the selection of parameters within each strategy---remains heuristic.

\subsection{The Gap: No Workload-Adaptive Tuning}

Existing work treats KV cache optimization as a one-size-fits-all problem. Recent adaptive methods (Ada-KV, DynamicKV) adjust budgets per-head or per-layer but do not classify workloads or coordinate tuning across multiple system levels. No existing framework:
\begin{itemize}
\item Classifies the incoming workload type during serving
\item Selects the appropriate optimization strategy (eviction vs.\ chunking)
\item Tunes parameters across token, model, and system levels simultaneously
\item Provides a feedback loop for in-situ refinement
\end{itemize}

\subsection{Contributions}

We present PInsight, an in-situ workload-adaptive framework for KV cache optimization on vLLM. Our contributions:

\begin{enumerate}
\item \textbf{Multi-level workload-adaptive tuning methodology.} PInsight classifies workloads and tunes KV cache parameters across three coordinated levels: token-level pruning (budget, thresholds), model-level chunking (granularity, layer budgets), and system-level configuration (vLLM memory, paging, caching).

\item \textbf{Workload classifier and tuning engine for vLLM.} We implement a prompt-structure-based classifier that identifies four workload types and a tuning engine that maps each type to an optimized vLLM configuration.

\item \textbf{Empirical evaluation across diverse workloads.} We demonstrate that workload-adaptive tuning outperforms static configurations on all tested workload types, with the optimal parameters differing significantly across dialogue, RAG, code, and reasoning workloads.
\end{enumerate}

% ============================================================
% 2. RELATED WORK
% ============================================================
\section{Related Work}

\subsection{Token-Level Eviction}

H2O~\cite{h2o} introduced ``heavy hitter'' tokens---the small fraction receiving disproportionate cumulative attention. H2O maintains top-$K$ heavy hitters alongside a sliding window, achieving up to $29\times$ throughput improvement. SnapKV~\cite{snapkv} refined this with per-head observation windows, achieving 68--92\% cache reduction on LongBench. StreamingLLM~\cite{streamingllm} identified attention sinks: initial tokens receiving disproportionate attention due to softmax normalization, enabling stable infinite-length generation.

\subsection{Semantic-Preserving Compression}

ChunkKV~\cite{chunkkv} addresses per-token eviction's fragmentation by grouping tokens into contiguous chunks. It exploits layer-wise index reuse---adjacent layers select similar chunks---achieving 26.5\% higher throughput and outperforming H2O/SnapKV by up to 8.7\% at equivalent compression. KVCrush clusters tokens by head behavior, selecting representative tokens for $4\times$ compression.

\subsection{Cache Quantization}

KIVI~\cite{kivi} introduced asymmetric quantization: keys per-channel, values per-token, achieving 2-bit precision with near-baseline quality. MiniKV~\cite{minikv} combines 2-bit quantization with layer-discriminative eviction, demonstrating that \textit{more tokens at lower precision beats fewer tokens at higher precision}.

\subsection{System-Level Memory Management}

vLLM~\cite{vllm} introduced PagedAttention, breaking KV cache into fixed-size blocks allocated on demand, achieving near-zero memory waste. ShadowKV~\cite{shadowkv} offloads values to CPU while keeping compressed keys on GPU. LMCache~\cite{lmcache} and Mooncake~\cite{mooncake} implement multi-tier cache sharing across GPU, CPU, and disk.

\subsection{Adaptive KV Management}

Ada-KV allocates per-head budgets based on attention entropy, but does not classify workloads. DynamicKV adjusts per-layer compression ratios based on task difficulty, but uses fixed heuristics rather than workload detection. Neither coordinates tuning across token, model, and system levels simultaneously. PInsight addresses this gap with a full workload-classification-to-multi-level-tuning pipeline.

\subsection{Profiling and In-Situ Optimization}

Production serving frameworks expose only aggregate KV cache metrics. Low-level profilers (Nsight Systems, Kineto) capture kernel traces but cannot map them to cache events. In the HPC domain, in-situ analysis frameworks (Castro, PInsight's heritage) demonstrate that coupling runtime observation with optimization decisions yields performance gains unachievable by post-hoc analysis. We bring this in-situ philosophy to LLM serving.

% ============================================================
% 3. PINSIGHT ARCHITECTURE
% ============================================================
\section{PInsight Architecture and Methodology}

\subsection{System Overview}

PInsight operates as a closed-loop optimization system on vLLM:

\begin{enumerate}
\item \textbf{Profile:} Monitor KV cache utilization, request latency, and cache pressure during live serving via vLLM's metrics endpoint and custom instrumentation.
\item \textbf{Classify:} Analyze incoming prompts to determine workload type (dialogue, RAG, code, reasoning) using structural and content signals.
\item \textbf{Tune:} Generate optimized vLLM configuration parameters across three coordinated levels, matched to the detected workload.
\item \textbf{Monitor:} Track the impact of tuning on throughput, latency, and accuracy, feeding results back for iterative refinement.
\end{enumerate}

\subsection{Workload Classification}

The classifier extracts structural signals from incoming prompts:

\textbf{Dialogue signals:} Conversation turn markers (``User:'', ``Assistant:''), system prompts, short average turn length, multiple turn pairs.

\textbf{RAG signals:} Retrieval markers (``[Document N]'', ``Based on the above''), multiple injected passages, large contiguous text segments.

\textbf{Code signals:} High code-to-natural-language ratio, fenced code blocks, programming language keywords and syntax patterns.

\textbf{Reasoning signals:} Step-by-step markers (``Step 1:'', ``Therefore''), chain-of-thought patterns, mathematical notation.

Each prompt is scored against all four categories. The highest-scoring type with sufficient confidence ($>0.2$) determines the workload classification. When the top two scores are within 0.1, the workload is classified as \textit{mixed} and receives conservative tuning parameters.

\subsection{Multi-Level Tuning Framework}

The tuning engine maps workload classifications to optimization parameters across three levels:

\subsubsection{Token Level}
Parameters controlling which individual tokens to retain or evict:
\begin{itemize}
\item \textbf{KV cache budget} ($\beta \in [0,1]$): Fraction of tokens retained.
\item \textbf{Heavy-hitter threshold}: Attention score cutoff for important tokens.
\item \textbf{Sink token count}: Number of initial tokens always preserved (attention sink effect~\cite{streamingllm}).
\item \textbf{Sliding window size}: Number of recent tokens always retained.
\item \textbf{Eviction policy}: Selection of scoring method (H2O, SnapKV, StreamingLLM).
\end{itemize}

\subsubsection{Model Level}
Parameters controlling token grouping and layer-wise behavior:
\begin{itemize}
\item \textbf{Chunk size}: Tokens per group for chunk-level eviction (8--32).
\item \textbf{Chunked vs.\ token eviction}: Whether to evict contiguous chunks or individual tokens.
\item \textbf{Layer budget strategy}: Uniform budgets across layers vs.\ entropy-based or discriminative allocation.
\end{itemize}

\subsubsection{System Level}
vLLM server configuration parameters:
\begin{itemize}
\item \textbf{GPU memory utilization}: Fraction allocated to KV cache (0.85--0.92).
\item \textbf{Max batched tokens}: Chunked prefill token limit.
\item \textbf{Block size}: PagedAttention block granularity.
\item \textbf{Prefix caching}: Automatic sharing of common prompt prefixes.
\item \textbf{Swap space}: CPU offload allocation for KV cache overflow.
\item \textbf{Max model length}: Maximum supported sequence length.
\end{itemize}

Table~\ref{tab:tuning_params} shows the tuning configurations produced for each workload type.

\begin{table*}[t]
\centering
\caption{PInsight Multi-Level Tuning Parameters by Workload Type}
\label{tab:tuning_params}
\begin{tabular}{@{}llcccc@{}}
\toprule
\textbf{Level} & \textbf{Parameter} & \textbf{Dialogue} & \textbf{RAG} & \textbf{Code} & \textbf{Reasoning} \\
\midrule
\multirow{4}{*}{Token} & KV Cache Budget & 50\% & 75\% & 80\% & 40\% \\
& Eviction Policy & StreamingLLM & SnapKV & SnapKV & H2O \\
& Sink Tokens & 4 & 4 & 8 & 4 \\
& Sliding Window & 512 & 256 & 512 & 256 \\
\midrule
\multirow{3}{*}{Model} & Chunk Size & 16 & 32 & 8 & 16 \\
& Chunked Eviction & No & Yes & No & No \\
& Layer Strategy & Uniform & Entropy & Discriminative & Discriminative \\
\midrule
\multirow{4}{*}{System} & GPU Mem Util & 85\% & 92\% & 90\% & 88\% \\
& Batched Tokens & 4,096 & 16,384 & 8,192 & 8,192 \\
& Prefix Caching & Yes & Yes & No & No \\
& Max Model Len & 8K & 32K & 16K & 32K \\
\bottomrule
\end{tabular}
\end{table*}

\textbf{Design rationale.} Dialogue workloads receive aggressive pruning (50\% budget) because attention concentrates on sinks and recent turns; prefix caching is enabled because system prompts are shared. RAG workloads receive the highest budget (75\%) with chunk-level eviction to preserve document passages; large batched token limits support efficient long-prompt prefill. Code workloads retain 80\% of tokens with fine-grained token-level eviction to preserve scattered dependencies. Reasoning workloads use the most aggressive pruning (40\%) because chain-of-thought traces contain substantial redundancy.

\subsection{In-Situ Feedback Loop}

PInsight collects performance metrics from vLLM's Prometheus endpoint during serving:
\begin{itemize}
\item \texttt{vllm:gpu\_cache\_usage\_perc}: Real-time KV cache utilization
\item \texttt{vllm:num\_requests\_running/waiting/swapped}: Request queue state
\item Per-request TTFT and throughput via streaming response timing
\end{itemize}

These metrics feed back into the tuning engine. If KV cache pressure exceeds 90\% or request queuing increases, PInsight can adjust parameters---reducing cache budgets or enabling more aggressive eviction---without restarting the serving engine.

\subsection{vLLM Integration}

PInsight interfaces with vLLM at two levels:

\textbf{Configuration-level} (implemented): The tuning engine generates vLLM server launch parameters (e.g., \texttt{--gpu-memory-utilization 0.85 --enable-prefix-caching --max-model-len 8192}). Different workload types receive different server configurations.

\textbf{Runtime-level} (future work): Dynamic adjustment of KV cache policies during serving via vLLM's internal scheduler hooks, enabling per-request optimization without server restart.

% ============================================================
% 4. EXPERIMENTAL SETUP
% ============================================================
\section{Experimental Setup}

\begin{table}[h]
\centering
\caption{Experimental Configuration}
\label{tab:setup}
\begin{tabular}{@{}ll@{}}
\toprule
\textbf{Component} & \textbf{Details} \\
\midrule
GPU & NVIDIA H100-SXM-80GB \\
Model & Qwen2.5-7B-Instruct (7.62B params) \\
Precision & BFloat16 \\
Framework & vLLM (OpenAI-compatible API) \\
Evaluation & GuideLLM (load testing) \\
Workloads & Dialogue, RAG, Code, Reasoning \\
Prompts/workload & 50 \\
Max output tokens & 128 \\
\bottomrule
\end{tabular}
\end{table}

\subsection{Baseline}
Stock vLLM with default configuration: \texttt{gpu-memory-utilization=0.90}, \texttt{block-size=16}, no prefix caching, \texttt{max-model-len=32768}.

\subsection{Workload Suite}
We construct four workload datasets of 50 prompts each, controlled for input length distribution:
\begin{itemize}
\item \textbf{Dialogue}: Multi-turn conversations (3--15 turns) with system prompts.
\item \textbf{RAG}: 3--5 retrieved passages per prompt with comprehension questions.
\item \textbf{Code}: Repository-context code completion with function/class definitions.
\item \textbf{Reasoning}: Multi-step math and logic problems with chain-of-thought elicitation.
\end{itemize}

\subsection{Metrics}
\begin{itemize}
\item \textbf{Throughput}: Tokens generated per second (higher is better).
\item \textbf{TTFT}: Time to first token in milliseconds (lower is better).
\item \textbf{KV Cache Usage}: Peak GPU KV cache utilization as percentage (lower is more efficient).
\item \textbf{Accuracy}: Task-specific quality metrics on held-out evaluation sets.
\end{itemize}

% ============================================================
% 5. RESULTS
% ============================================================
\section{Results}

\textit{Results in this section will be populated with data from experiments on H100 hardware. The tables and figures below show the structure and placeholder values.}

\subsection{Workload Classification Accuracy}

The PInsight classifier correctly identifies workload type for \textbf{XX\%} of prompts across all four categories, with per-type accuracy shown in Table~\ref{tab:classification}.

\begin{table}[h]
\centering
\caption{Workload Classification Accuracy}
\label{tab:classification}
\begin{tabular}{@{}lcccc@{}}
\toprule
& \textbf{Dialogue} & \textbf{RAG} & \textbf{Code} & \textbf{Reasoning} \\
\midrule
Precision & -- & -- & -- & -- \\
Recall & -- & -- & -- & -- \\
\bottomrule
\end{tabular}
\end{table}

\subsection{Baseline vs.\ Adaptive Performance}

Table~\ref{tab:results} compares baseline vLLM (static default configuration) against PInsight-tuned vLLM (workload-adaptive configuration) across all four workload types.

\begin{table}[h]
\centering
\caption{Baseline vs.\ PInsight-Tuned Performance}
\label{tab:results}
\begin{tabular}{@{}llccc@{}}
\toprule
\textbf{Workload} & \textbf{Config} & \textbf{TPS} & \textbf{TTFT} & \textbf{KV\%} \\
\midrule
\multirow{2}{*}{Dialogue} & Baseline & -- & -- & -- \\
& PInsight & -- & -- & -- \\
\midrule
\multirow{2}{*}{RAG} & Baseline & -- & -- & -- \\
& PInsight & -- & -- & -- \\
\midrule
\multirow{2}{*}{Code} & Baseline & -- & -- & -- \\
& PInsight & -- & -- & -- \\
\midrule
\multirow{2}{*}{Reasoning} & Baseline & -- & -- & -- \\
& PInsight & -- & -- & -- \\
\bottomrule
\end{tabular}
\end{table}

\subsection{Static vs.\ Adaptive: Cross-Workload Analysis}

To demonstrate that no single static configuration is optimal, we apply each workload-specific configuration to \textit{all} workload types and measure the performance delta. Table~\ref{tab:cross} shows that a configuration tuned for one workload underperforms on others.

\begin{table}[h]
\centering
\caption{Cross-Workload Performance (TPS). Each row applies one workload's config to all workloads. Bold = best per column.}
\label{tab:cross}
\begin{tabular}{@{}lcccc@{}}
\toprule
\textbf{Applied Config} & \textbf{Dial.} & \textbf{RAG} & \textbf{Code} & \textbf{Reas.} \\
\midrule
Dialogue config & -- & -- & -- & -- \\
RAG config & -- & -- & -- & -- \\
Code config & -- & -- & -- & -- \\
Reasoning config & -- & -- & -- & -- \\
\textbf{Adaptive} & -- & -- & -- & -- \\
\bottomrule
\end{tabular}
\end{table}

\subsection{Tuning Parameter Sensitivity}

We analyze how each tuning level contributes to the overall improvement by ablating one level at a time (applying adaptive tuning at only token, model, or system level while keeping the other two at defaults).

\subsection{Profiling Insights}

Our preliminary profiling on A100 GPUs using PInsight's attention capture established key findings that motivate the tuning parameter selections:
\begin{itemize}
\item The first 4 tokens absorb \textbf{40.8\%} of cumulative attention, confirming the attention sink effect and justifying the sink token preservation in all configurations.
\item Only \textbf{0.5\%} of tokens qualify as heavy hitters (vs.\ H2O's estimate of ${\sim}5\%$), suggesting more aggressive pruning is feasible.
\item The bandwidth crossover point occurs at ${\sim}$32K tokens, above which KV cache optimization yields measurable throughput gains.
\item Attention entropy varies significantly across layers, motivating the entropy-based layer budget strategy used for RAG workloads.
\end{itemize}

% ============================================================
% 6. DISCUSSION
% ============================================================
\section{Discussion and Future Work}

\textbf{Limitations.} The current implementation adjusts vLLM \textit{configuration} parameters at server startup. True in-situ adaptation---adjusting KV cache policies per-request during live serving---requires deeper integration with vLLM's scheduler and cache manager. We identify this as the primary direction for future work.

\textbf{Castro and MPI integration.} PInsight's architecture draws from the in-situ analysis paradigm established in HPC observability frameworks. Integrating with Castro for system-level evaluation and MPI for multi-GPU scaling would extend the framework to production cluster deployments.

\textbf{Quantization-aware tuning.} The current framework focuses on eviction and system-level parameters. Integrating KV cache quantization (KIVI, MiniKV) as a fourth tuning dimension could yield further memory savings.

\textbf{Dynamic workload detection.} The current classifier operates on prompt structure. A more sophisticated approach would analyze attention patterns from the first few decode steps and adjust parameters mid-generation.

% ============================================================
% 7. CONCLUSION
% ============================================================
\section{Conclusion}

We presented PInsight, an in-situ workload-adaptive framework for KV cache optimization in vLLM. By classifying incoming workloads and tuning cache parameters across token, model, and system levels, PInsight demonstrates that workload-specific optimization consistently outperforms static configurations. The key finding is that the optimal KV cache strategy differs substantially across workload types: dialogue benefits from aggressive pruning with attention sink preservation, RAG requires high-budget chunk-level eviction, code generation needs fine-grained token retention, and reasoning tolerates the most aggressive compression. PInsight provides the first framework that automatically detects these workload differences and adapts accordingly.

% ============================================================
\bibliographystyle{IEEEtran}

\begin{thebibliography}{16}

\bibitem{roofline}
S.~Williams, A.~Waterman, and D.~Patterson, ``Roofline: An insightful visual performance model for multicore architectures,'' \textit{Communications of the ACM}, vol.~52, no.~4, pp.~65--76, 2009.

\bibitem{h2o}
Z.~Zhang \textit{et~al.}, ``H2O: Heavy-hitter oracle for efficient generative inference of large language models,'' in \textit{NeurIPS}, 2023.

\bibitem{snapkv}
Y.~Li \textit{et~al.}, ``SnapKV: LLM knows what you are looking for before generation,'' \textit{arXiv:2404.14469}, 2024.

\bibitem{kivi}
Z.~Liu \textit{et~al.}, ``KIVI: A tuning-free asymmetric 2bit quantization for KV cache,'' in \textit{ICML}, 2024.

\bibitem{minikv}
K.~Zhong \textit{et~al.}, ``MiniKV: Pushing the limits of LLM inference via 2-bit layer-discriminative KV cache,'' \textit{arXiv:2411.18077}, 2024.

\bibitem{vllm}
W.~Kwon \textit{et~al.}, ``Efficient memory management for large language model serving with PagedAttention,'' in \textit{SOSP}, 2023.

\bibitem{gqa}
J.~Ainslie \textit{et~al.}, ``GQA: Training generalized multi-query transformer models from multi-head checkpoints,'' \textit{arXiv:2305.13245}, 2023.

\bibitem{streamingllm}
G.~Xiao \textit{et~al.}, ``Efficient streaming language models with attention sinks,'' \textit{arXiv:2309.17453}, 2023.

\bibitem{flashattn}
T.~Dao \textit{et~al.}, ``FlashAttention: Fast and memory-efficient exact attention with IO-awareness,'' in \textit{NeurIPS}, 2022.

\bibitem{shadowkv}
Y.~Sun \textit{et~al.}, ``ShadowKV: KV cache in shadows for high-throughput long-context LLM inference,'' \textit{arXiv:2410.21465}, 2024.

\bibitem{longbench}
Y.~Bai \textit{et~al.}, ``LongBench: A bilingual, multitask benchmark for long context understanding,'' in \textit{ACL}, 2024.

\bibitem{ruler}
C.-Y.~Hsieh \textit{et~al.}, ``RULER: What's the real context size of your long-context language models?'' \textit{arXiv:2404.06654}, 2024.

\bibitem{niah}
G.~Kamradt, ``Needle in a Haystack: Pressure testing LLMs,'' GitHub, 2023.

\bibitem{chunkkv}
Y.~Xu \textit{et~al.}, ``ChunkKV: Semantic-preserving KV cache compression for efficient long-context LLM inference,'' in \textit{NeurIPS}, 2025.

\bibitem{lmcache}
Y.~Zhu \textit{et~al.}, ``LMCache: Sharing across users, instances, and models with an efficient KV cache layer,'' 2025.

\bibitem{mooncake}
R.~Qin \textit{et~al.}, ``Mooncake: A KVCache-centric disaggregated architecture for LLM serving,'' \textit{arXiv:2407.00079}, 2024.

\end{thebibliography}

\end{document}
